% Generated draft: requires mathematical and source review.
\documentclass[11pt]{article}
\usepackage[margin=25mm]{geometry}
\usepackage{fontspec}
\setmainfont{Latin Modern Roman}
\usepackage{amsmath,amssymb,mathtools,unicode-math}
\setmathfont{Latin Modern Math}
\usepackage{enumitem,needspace,xurl,hyperref}
\setlength{\parindent}{0pt}
\setlength{\parskip}{5pt}
\setlength{\emergencystretch}{3em}
\newcommand{\N}{\mathbb{N}}
\newcommand{\Z}{\mathbb{Z}}
\newcommand{\Q}{\mathbb{Q}}
\newcommand{\R}{\mathbb{R}}
\newcommand{\C}{\mathbb{C}}
\newcommand{\F}{\mathbb{F}}
\newcommand{\A}{\mathbb{A}}
\newcommand{\PP}{\mathbb{P}}
\newcommand{\E}{\mathbb{E}}
\newcommand{\eps}{\varepsilon}
\newcommand{\dd}{\,\mathrm{d}}
\begin{document}
\section*{UnsolvedMath ID: 30001614}
\textbf{Source problem code:} catalogue entry
\subsection{Definitions and mathematical statement}
The supplied record concerns two questions about spaces denoted by
\[
K^p_\theta=H^p\cap\theta H^p_0,
\]
where \(\theta\) is an inner function. This formula is reproduced exactly as supplied. The underlying Hardy-space domain, the admissible range of \(p\), and the definition of \(H^p_0\) are not specified.

There is a material notation issue. Under the usual unit-disc convention, \(H^p\) denotes a Hardy space of holomorphic functions on \(\mathbb D=\{z\in\C:|z|<1\}\), \(H^p_0=\{f\in H^p:f(0)=0\}\), and an inner function is a bounded holomorphic function whose radial boundary values have modulus one almost everywhere on \(\partial\mathbb D\). With these conventions, the displayed formula gives \(K^p_\theta=\theta H^p_0\). The customary star-invariant-space formula instead uses conjugated boundary functions:
\[
H^p\cap\theta\overline{H^p_0},
\qquad
\overline{H^p_0}=\{\overline{u}:u\in H^p_0\},
\]
with the intersection interpreted through boundary values. These conventions explain the possible discrepancy; they are not a verified correction of the supplied record. The intended definition must be resolved before either question has an unambiguous target.

For any chosen interpretation of \(K^\infty_\theta\) and its inherited Hardy-space norm, let
\[
B_\theta=\{f\in K^\infty_\theta:\|f\|_\infty\leq1\}.
\]
A point \(f\in B_\theta\) is an extreme point if, whenever \(u,v\in B_\theta\) and \(0<t<1\) satisfy \(f=tu+(1-t)v\), one has \(u=v=f\). The first question is to characterize all extreme points of \(B_\theta\) for an arbitrary inner function \(\theta\).

The second question concerns solutions of
\[
f^n+g^n=h^n,
\qquad f,g,h\in K^p_\theta,\qquad n\in\Z,\quad n\geq3,
\]
where the equality is an identity of the represented functions. The source exempts triples spanning a one-dimensional subspace. Under the interpretation that the question asks whether every solution is proportional, its proposed assertion is
\[
\forall f,g,h\in K^p_\theta,\qquad
f^n+g^n=h^n
\ \Longrightarrow\ 
\dim_{\C}\operatorname{span}_{\C}\{f,g,h\}\leq1,
\]
for every inner \(\theta\), every admissible \(p\), and every integer \(n\geq3\). Equivalently, the question asks whether a solution spanning a subspace of dimension at least two can exist. This is a conditional rendering of the apparent nonproportional-solution question: the record does not explicitly state the nonexistence assertion, specify whether individual functions must be nonzero, or address the zero-dimensional span of the all-zero triple. No range of \(p\) is supplied, and membership of the powers in \(K^p_\theta\) is not asserted.
\subsection{Short English statement}
The record asks for a description of the extreme points of the unit ball of \(K^\infty_\theta\), and whether \(f^n+g^n=h^n\), for integers \(n\geq3\), can have solutions in \(K^p_\theta\) that are not all proportional. The space definition may be missing complex conjugation; the Hardy-space domain, range of \(p\), and treatment of zero functions also require clarification. This is a statement-only draft requiring human mathematical review.
\subsection{Sources}
\begin{itemize}
\item Local UnsolvedMath registry, supplied complete record: permanent ID \(30001614\), exact source problem code \texttt{catalogue entry}, title \emph{Extreme Points and Fermat Equations in Star-Invariant Spaces}, collection \emph{Oberwolfach Reports Open Problems}. The record gives proposed year 2010 and no named proposer in its proposer field; its background and literature assessment attribute both questions to Dyakonov.
\item Supplied source citation: \emph{Operator Theory and Harmonic Analysis (2011)}, Oberwolfach Reports, pp.\ 2813--2875, DOI \texttt{10.4171/owr/2010/49}. Original problem report: \url{https://doi.org/10.4171/owr/2010/49}. The background describes Oberwolfach Reports as documenting mathematical workshops and their open-problem discussions. The citation's 2011 label and the record's 2010 proposed year are retained without resolving their relationship.
\item Supplied literature source, labelled ``Dyakonov two-problem paper'': \url{https://arxiv.org/abs/1401.8031}.
\item Supplied literature source, labelled ``Dyakonov later extreme-point questions'': \url{https://doi.org/10.1007/s00020-023-02733-8}.
\item Supplied literature assessment checked at \texttt{2026-08-21}: the record labels the problem open, attributes both questions for general star-invariant subspaces to Dyakonov, and reports that related later extreme-point work was found but no general solution of the extracted pair was found. This is supplied dated curation derived from the OWR material, not an independently verified assessment. The record requests independent verification of time-sensitive status and cited claims.
\item The supplied statement status is \texttt{corrected\_verified}; its verification note reports a manual semantic, typographic, and LaTeX review claiming a self-contained, meaning-preserving cleanup. That supplied label does not resolve the formulation issues identified here or constitute independent mathematical verification.
\end{itemize}
\end{document}
